% =====================================================================
% tesi/capitoli/29-assi-della-collezione.tex
% =====================================================================
% copre: pokedex-home-completo/STUDIO-05-gli-assi-che-non-contavamo.md
% copre: pokedex-home-completo/STUDIO-05-gli-assi-che-non-contavamo.md#studio-05-gli-assi-che-non-contavamo-e-una-misura-sulle-mosse-perdute
% copre: pokedex-home-completo/STUDIO-05-gli-assi-che-non-contavamo.md#perché-questo-studio-esiste
% copre: pokedex-home-completo/STUDIO-05-gli-assi-che-non-contavamo.md#le-mosse-che-sui-titoli-per-console-corrente-non-esistono-più
% copre: pokedex-home-completo/STUDIO-05-gli-assi-che-non-contavamo.md#i-fiocchi-che-soltanto-un-gioco-vecchio-conferisce
% copre: pokedex-home-completo/STUDIO-05-gli-assi-che-non-contavamo.md#le-combinazioni-di-sfera-che-soltanto-un-gioco-vecchio-produce
% copre: pokedex-home-completo/STUDIO-05-gli-assi-che-non-contavamo.md#un-quarto-asse-che-è-di-natura-diversa
% copre: pokedex-home-completo/STUDIO-05-gli-assi-che-non-contavamo.md#la-misura-quante-delle-mosse-perdute-i-nostri-lotti-già-portano
% copre: pokedex-home-completo/STUDIO-05-gli-assi-che-non-contavamo.md#due-conferme-indipendenti-di-decisioni-già-prese
% copre: pokedex-home-completo/STUDIO-05-gli-assi-che-non-contavamo.md#una-seconda-testimonianza-sulla-via-del-canale-ricostruito
% copre: pokedex-home-completo/STUDIO-05-gli-assi-che-non-contavamo.md#che-cosa-entra-nel-lavoro-e-in-quale-ordine
% copre: pokedex-home-completo/STUDIO-06-le-enumerazioni-trasversali.md
% copre: pokedex-home-completo/STUDIO-06-le-enumerazioni-trasversali.md#studio-06-le-enumerazioni-trasversali-e-la-prima-misura-che-concorda-su-tutta-la-riga
% copre: pokedex-home-completo/STUDIO-06-le-enumerazioni-trasversali.md#che-cosa-è-stato-letto-e-che-cosa-no
% copre: pokedex-home-completo/STUDIO-06-le-enumerazioni-trasversali.md#l-asse-del-sesso-è-chiuso-e-le-due-misure-indipendenti-concordano-esattamente
% copre: pokedex-home-completo/STUDIO-06-le-enumerazioni-trasversali.md#le-sfide-del-deposito-e-la-fonte-primaria-che-avevamo-in-casa-senza-guardarla
% copre: pokedex-home-completo/STUDIO-06-le-enumerazioni-trasversali.md#i-marchi-di-provenienza-e-un-fatto-che-cambia-come-si-enuncia-la-loro-permanenza
% copre: pokedex-home-completo/STUDIO-06-le-enumerazioni-trasversali.md#due-enumerazioni-nuove-e-che-cosa-aggiungono-davvero
% copre: pokedex-home-completo/STUDIO-06-le-enumerazioni-trasversali.md#il-controllo-all-indietro-che-adr-049-impone-e-i-candidati-che-questo-cluster-porta
% copre: pokedex-home-completo/STUDIO-06-le-enumerazioni-trasversali.md#i-doni-del-deposito-stesso-che-nessuna-nostra-enumerazione-conteneva
% copre: pokedex-home-completo/STUDIO-06-le-enumerazioni-trasversali.md#i-sottolivellati-e-i-cromatici-garantiti-cioè-due-assi-con-la-loro-fonte
% copre: pokedex-home-completo/STUDIO-06-le-enumerazioni-trasversali.md#che-cosa-resta-aperto-e-in-quale-ordine
% copre: pokedex-home-completo/STUDIO-07-scambi-in-gioco-e-incontri-condizionati.md
% copre: pokedex-home-completo/STUDIO-07-scambi-in-gioco-e-incontri-condizionati.md#studio-07-gli-scambi-in-gioco-e-gli-incontri-che-una-condizione-sblocca
% copre: pokedex-home-completo/STUDIO-07-scambi-in-gioco-e-incontri-condizionati.md#che-cosa-unisce-le-due-classi-e-perché-non-è-la-rarità
% copre: pokedex-home-completo/STUDIO-07-scambi-in-gioco-e-incontri-condizionati.md#gli-scambi-in-gioco-contati-sulla-fonte-di-primo-livello
% copre: pokedex-home-completo/STUDIO-07-scambi-in-gioco-e-incontri-condizionati.md#gli-incontri-condizionati-e-una-correzione-all-esempio-da-cui-il-lavoro-è-nato
% copre: pokedex-home-completo/STUDIO-07-scambi-in-gioco-e-incontri-condizionati.md#che-cosa-dice-adr-049-su-queste-due-classi-e-la-domanda-che-resta-all-utente
% copre: pokedex-home-completo/STUDIO-07-scambi-in-gioco-e-incontri-condizionati.md#che-cosa-resta-da-fare-su-queste-due-classi
% copre: pokedex-home-completo/MOSSE-PERDUTE.md
% copre: pokedex-home-completo/MOSSE-PERDUTE.md#le-mosse-perdute-derivate-dai-dati-e-non-citate
% copre: pokedex-home-completo/MOSSE-PERDUTE.md#il-confronto-con-la-testimonianza
% copre: pokedex-home-completo/MOSSE-PERDUTE.md#che-cosa-i-nostri-lotti-gia-portano
% copre: pokedex-home-completo/FIOCCHI.md
% copre: pokedex-home-completo/FIOCCHI.md#l-asse-dei-fiocchi-enumerazione-e-copertura
% copre: pokedex-home-completo/FIOCCHI.md#l-asse-intero-e-la-parte-che-i-nostri-formati-sanno-rappresentare
% copre: pokedex-home-completo/DIFFERENZE-DI-SESSO.md
% copre: pokedex-home-completo/DIFFERENZE-DI-SESSO.md#le-specie-con-differenze-di-sesso-visibili
% copre: pokedex-home-completo/CENSIMENTO-SCAMBI.md
% copre: pokedex-home-completo/CENSIMENTO-SCAMBI.md#censimento-degli-scambi-in-gioco-tutte-le-generazioni
% copre: pokedex-home-completo/CENSIMENTO-CONDIZIONATI.md
% copre: pokedex-home-completo/CENSIMENTO-CONDIZIONATI.md#censimento-degli-incontri-che-una-condizione-sblocca
% copre: pokedex-home-completo/CENSIMENTO-CONDIZIONATI.md#quadro-d-insieme
% copre: pokedex-home-completo/CENSIMENTO-CONDIZIONATI.md#le-specie-che-fra-le-vie-selvatiche-vengono-solo-da-una-condizione
% copre: pokedex-home-completo/CENSIMENTO-CONDIZIONATI.md#le-aree-che-ospitano-una-specie-sola-e-non-compare-altrove-nel-titolo
% copre: pokedex-home-completo/CENSIMENTO-CONDIZIONATI.md#i-luoghi-dedicati-a-una-sola-specie
% copre: pokedex-home-completo/CONFRONTO-FOGLIO-LIVINGDEX.md
% copre: pokedex-home-completo/CONFRONTO-FOGLIO-LIVINGDEX.md#confronto-fra-la-nostra-enumerazione-e-il-foglio-comunitario-del-living-dex
% copre: pokedex-home-completo/CONFRONTO-FOGLIO-LIVINGDEX.md#il-conto-delle-due-enumerazioni
% copre: pokedex-home-completo/CONFRONTO-LIVINGDEX-POKEPC.md
% copre: pokedex-home-completo/CONFRONTO-LIVINGDEX-POKEPC.md#confronto-fra-la-nostra-enumerazione-e-quella-di-pokepc-classic
% copre: pokedex-home-completo/CONFRONTO-LIVINGDEX-POKEPC.md#il-conto-delle-due-enumerazioni
% copre: pokedex-home-completo/CONFRONTO-LIVINGDEX-POKEPC.md#che-cosa-pokepc-colloca-in-una-casella-per-classe
% copre: pokedex-home-completo/CONFRONTO-LIVINGDEX-POKEPC.md#le-disposizioni-sorelle-che-non-concordano-fra-loro
% copre: pokedex-home-completo/EVENTI-SENZA-CARTA.md
% copre: pokedex-home-completo/EVENTI-SENZA-CARTA.md#distribuzioni-senza-carta-la-coda-di-produzione
% copre: pokedex-home-completo/EVENTI-SENZA-CARTA.md#i-due-casi-in-cui-un-campo-non-è-un-valore
% copre: pokedex-home-completo/ID-NOTEVOLI.md
% copre: pokedex-home-completo/ID-NOTEVOLI.md#gli-identificativi-di-allenatore-notevoli-e-quanto-ne-copriamo
% verificato-al-commit: 6b2dc05
% ---------------------------------------------------------------------

\chapter{Gli assi di una collezione, e come si enumerano}
\label{cap:assi}

Il capitolo~\ref{cap:enumerazione} ha stabilito che cosa si possa ancora ottenere e da quali titoli, cioè ha misurato l'\emph{ottenibilità}. Questo capitolo tratta una domanda che gli sta prima e che il lavoro ha impiegato settimane a formulare correttamente: che cosa esattamente si stia contando. La differenza non è filosofica. Un'enumerazione sbagliata non produce alcun errore visibile, perché ogni conto torna rispetto a se stesso; produce una collezione che si dichiara completa mentre è vuota su intere dimensioni che nessuno aveva pensato di contare.

\section{Tre assi diventano sei}

Fino alla lettura sistematica delle fonti di comunità, il lavoro contava tre assi. Le specie, che sono milleventicinque. Le voci di forma, che dopo le decisioni sul perimetro sono trecentoquarantadue. Gli esemplari da distribuzione, che sono duemilaseicentottantasei sotto scadenza. Su quei tre le enumerazioni esistevano e i generatori producevano.

La lettura del raccoglitore comunitario su ciò che si perde con la chiusura della banca \cite{bank-exclusives} ha mostrato che gli assi sono almeno sei, e che i tre nuovi non sono raffinamenti dei primi ma dimensioni indipendenti nel senso preciso che una collezione completa sui primi tre può essere vuota sui secondi. Il motivo per cui erano sfuggiti è strutturale: ciascuno dei tre nuovi non è una proprietà di \emph{quali} esemplari si possiedano, ma di quali proprietà quegli esemplari portino addosso.

Il primo è quello delle mosse rimosse dai titoli per console corrente. Chi le vuole nel deposito deve trasferire un esemplare che le conosca, perché non esiste modo di insegnarle di nuovo. È indipendente dagli altri perché una specie può essere presente e la sua mossa perduta, e ha una proprietà che gli altri non hanno, cioè che una mossa viaggia su qualunque esemplare che la conosca: non richiede un esemplare in più, richiede la scelta giusta di quale esemplare trasferire.

Il secondo è quello dei fiocchi, che sono contrassegni conferiti da un gioco e conservati dal dato attraverso i trasferimenti. Ogni generazione ne ha conferiti che le successive non conferiscono, e due di essi si ottengono soltanto purificando esemplari particolari nei due giochi per console fissa della terza generazione, che il lavoro non aveva mai considerato una sorgente.

Il terzo è quello delle combinazioni fra specie e sfera di cattura che nessun gioco moderno può produrre. Condivide la forma con l'asse dei cromatici, cioè moltiplica una specie per una proprietà, ed entrambi vanno misurati e non moltiplicati: il numero di combinazioni concepibili non è il numero di combinazioni esistenti, e assumere il primo per il secondo è il difetto di metodo che questo lavoro ha registrato con una decisione dedicata.

A questi si aggiunge una quarta dimensione che è di natura diversa e va tenuta separata, cioè le \emph{sfide} interne al servizio di deposito. Non sono un asse di collezione ma di completamento, e la ragione pratica per cui non vanno mescolate è che una sfida si soddisfa con un esemplare che spesso serve già a un altro asse: non aggiunge caselle, vincola la provenienza di caselle che si sarebbero riempite comunque.

\section{Una misura invece di una lettura}

Sull'asse delle mosse il lavoro ha adottato un metodo che vale oltre il caso e che questo capitolo enuncia come regola: una fonte di quinto livello che elenca cose perdute non si legge per crederle, si trasforma in un predicato e lo si esegue sui propri dati.

Applicato alla lettera, il predicato ha dato che ventidue delle sessantatre mosse elencate dalla testimonianza erano già dentro i lotti prodotti e quarantuno no. Il risultato non dipende dall'affidabilità della fonte per la parte che conta, cioè quali mosse i lotti contengano, e dipende da essa soltanto per quali siano da cercare: se la testimonianza avesse elencato una mossa in più o in meno, la misura resterebbe corretta su tutte le altre.

Il passo successivo ha sostituito la testimonianza con una derivazione dai dati del verificatore, e il confronto fra le due liste è l'esito più utile dei due. Una mossa si dice perduta quando è resa inefficace in tutti i titoli per console corrente insieme, cioè presente nei dati e non utilizzabile; dall'insieme si tolgono le mosse che nessun esemplare può conoscere, perché non sono perdute ma non sono mai state possedute. Intersecando i cinque contesti che portano una tabella di mosse rese inefficaci si ottengono ottanta mosse; la testimonianza ne elencava sessantatre e in comune sono quarantacinque.

Lo scarto non è un disaccordo ed è precisamente il genere di cosa che una lettura non avrebbe scoperto. I cinque contesti intersecati sono i soli che portino quella tabella, e le riedizioni della prima generazione per console corrente non ne portano alcuna, quindi su di esse la derivazione non dice nulla; la testimonianza invece le considera e dichiara che le mosse della prima generazione si recuperano di là. Le due liste concordano una volta tenuto conto di quel titolo. Al netto delle trentacinque voci recuperabili restano quarantacinque mosse davvero perdute, di cui diciassette già dentro i lotti prodotti e ventotto no: quest'ultimo numero definisce un lotto di lavoro piccolo e chiuso, che non richiede altra ricerca perché la lista è esattamente quella.

\section{L'asse del sesso, e la prima enumerazione che chiude in entrambi i versi}

Le differenze di sesso visibili sono un asse che l'utente aveva deciso di contare e per il quale mancava la fonte. La pagina enciclopedica dedicata \cite{bulbapedia-differenze-sesso} si è rivelata enumerabile a macchina perché ogni riga delle sue tabelle porta davanti il numero di catalogo a quattro cifre, e l'estrazione dà centodue specie contro le centodue che la pagina stessa dichiara in prosa: la fonte concorda con se stessa su tutte e otto le generazioni interessate.

Il confronto con il foglio comunitario del catalogo vivente, che è una seconda misura nata da una fonte che non conosce la prima, dà zero divergenze in entrambi i versi. È la prima volta che una enumerazione di questo lavoro chiude con scarto nullo su entrambi i lati invece che su uno solo, ed è la ragione per cui il criterio di chiusura adottato qui prescrive di misurare due numeri in versi opposti invece di uno: una enumerazione non si dice chiusa in assoluto, si dice chiusa rispetto a una fonte, e la misura in un verso solo nasconde metà del disaccordo.

Ne è venuta anche una correzione a un numero che il lavoro portava, cioè centotre specie invece di centodue. Le righe del foglio erano duecentosei e sono giuste; il centotre nasceva dal dividerle per due assumendo che ogni specie valesse una coppia, mentre una specie ne porta quattro perché compare in due forme, ciascuna nei due sessi. È la medesima famiglia di difetto già registrata sull'asse dei fiocchi, cioè un denominatore assunto invece che misurato, e la sua firma è sempre la stessa: il conto torna e il denominatore è sbagliato.

\section{L'asse dei fiocchi, e la differenza fra rappresentabile e ottenibile}

L'enumerazione dei fiocchi ha prodotto lo stesso difetto in forma diversa, e vale mostrarlo perché è il caso in cui la tentazione di enumerare dallo strumento sbagliato è più forte. Il formato della quarta e quinta generazione dichiara settantasei fiocchi distinti in dieci byte dell'esemplare, e la terza ne tiene trentadue in una parola sola: chi enumeri l'asse dal formato conclude che i fiocchi siano settantasei. La tabella dei nomi del verificatore ne elenca invece centosessantaquattro, e quello è l'asse; i settantasei sono la parte che i formati dei nostri lotti sanno rappresentare, cioè quelli esistiti fino alla quinta generazione.

Sulla parte rappresentabile la misura dice che sessantanove fiocchi su settantasei non sono portati da alcun esemplare prodotto, e la classificazione per via di conferimento dice che tutti e sessantanove appartengono a insiemi di generazione dalla terza alla settima, quindi si perdono con la chiusura. Il numero è dichiarato come limite superiore e non come valore esatto, per una ragione che il programma stesso registra: le riedizioni della quarta generazione per console corrente riconferiscono una parte di quei fiocchi, e il verificatore lo esprime dentro condizioni che il programma non interpreta. Distinguerle richiede di leggere quelle condizioni una per una, ed è lavoro dichiarato e non fatto.

\section{Le sfide del deposito, e una gerarchia di fonti che si rovescia}

L'asse delle sfide è il caso in cui il lavoro ha scoperto di avere in casa la fonte primaria senza guardarla. Le fonti erano tre e con coperture diverse: un post di comunità che elenca trentasette sfide esplicite più cinque invisibili \cite{sfide-impossibili}, una cartella di calcolo che un lettore di quel post ne ha derivato, e l'applicazione stessa, di cui l'utente aveva consegnato ventuno fotografie.

La lettura delle fotografie stabilisce che l'elenco dell'applicazione è di un ordine di grandezza più grande di quello che le due fonti derivate raccontano, perché quelle enumerano le sole sfide che la chiusura rende impossibili mentre l'applicazione le porta tutte: le sfide per ciascuno dei diciotto tipi, quelle per ciascuna delle venticinque indoli, quelle per il deposito in ciascuna sfera, il riempimento di ciascun catalogo regionale titolo per titolo, le sfide sul conto delle forme di una specie, e infine quelle di provenienza, che sono le sole sotto scadenza.

Le correzioni che la fonte primaria porta alle due derivate sono quattro, e tre alleggeriscono la scadenza. Due sfide che chiedono esemplari dai titoli di Kanto non sono sotto scadenza, perché quelle riedizioni si collegano al deposito da ottobre. Una sfida esiste nell'applicazione e non nell'archivio da cui il post discende, il che stabilisce che su questo asse la fonte primaria è l'applicazione e l'archivio è incompleto. Una terza sfida non è impossibile perché la specie richiesta è consegnata anche da un altro titolo. La quarta non alleggerisce nulla ed è la più utile per pianificare: una sfida di provenienza si soddisfa anche con un esemplare ricevuto in scambio da chi lo ha trasferito, quindi la scadenza vincola l'esistenza di quegli esemplari nel mondo e non la facoltà di registrarli.

Va dichiarato il limite della copertura, perché una copertura parziale non dichiarata somiglia a una completa: le fotografie riguardano la sola scheda principale, mentre le due schede accanto non sono state fotografate e le loro sfide restano non enumerate.

\section{Il controllo all'indietro, e le classi la cui via si è chiusa}

Il criterio di produzione adottato in questo lavoro dice che si produce un esemplare quando la sua sola via di provenienza non esiste più. La sua parte scomoda è che si applica anche all'indietro: ogni classe scartata va riletta chiedendosi se la sua via si sia chiusa nel frattempo, e senza quel controllo il criterio diventa una giustificazione a posteriori delle classi che si erano già scelte.

Le due enumerazioni comunitarie lette insieme forniscono esattamente la lista dei candidati. La prima è una cartella di ottocentouno voci fra scambi in gioco, doni, uova ed esemplari con cui si interagisce \cite{foglio-scambi-doni}, le cui due schede piccole, sui giochi derivati e sui giochi di deposito, sono una lista di vie chiuse. La seconda è una cartella di millequattrocentocinquantotto esemplari da evento interno più duecentotrentotto scambi \cite{foglio-eventi-interni}, la cui distribuzione per titolo dà novecentottantasei voci da titoli la cui via passa dalla banca e quattrocentosettantadue da titoli a via diretta.

Il secondo numero va letto con il criterio accanto, altrimenti conduce alla conclusione opposta al vero: quelle novecentottantasei voci sono in grandissima parte incontri ordinari, la cui via è aperta perché le cartucce sono possedute, e appartengono quindi alla lista di ciò che si ottiene giocando e non alla coda di produzione.

Le classi a via effettivamente chiusa che ne risultano sono undici, e comprendono i premi dei tornei per console fissa, i dischi allegati alle riviste, le ricompense di un titolo per console fissa distribuite attraverso il suo servizio in rete, i doni di un gioco derivato, le uova consegnate a soglie di deposito dal gestore per console fissa della terza generazione, gli scambi di un'applicazione del negozio in rete oggi dismesso, i dischi dimostrativi dei negozi e il sondatore dei sogni. Su quest'ultimo le due fonti danno numeri diversi senza contraddirsi, e la ragione va registrata perché si ripresenta ogni volta che si confrontano due enumerazioni: le tabelle del verificatore ne portano centottantuno perché contano ogni combinazione fra specie e livello, la fonte comunitaria ne elenca ventisei perché conta le specie e le forme. Il conto da usare per la coda è quello del verificatore, e l'altro serve da controllo di copertura.

Una riga di quella tabella va guardata prima delle altre perché tocca lavoro già fatto, ed è quella dei premi dei tornei: il lotto delle prime due generazioni li contiene e il verificatore li ha dichiarati conformi, mentre la fonte comunitaria afferma che almeno una di quelle voci non è trasferibile. Le due affermazioni non si escludono, perché la conformità di un esemplare e la sua trasferibilità sono controlli diversi fatti da soggetti diversi, ed è la verifica che il capitolo~\ref{cap:catena} ha poi chiuso.

\section{I doni del deposito stesso, e due assi con la loro fonte}

La lettura ha portato una classe che nessuna enumerazione del lavoro conteneva, cioè gli esemplari che il deposito consegna da sé: quello del primo avvio, la scelta fra i tre iniziali di Kanto, quello che si ottiene completando una sfida, i due legati ai sistemi di scambio, i terzetti che si ottengono depositando da ciascun titolo a via diretta, quello del primo trasferimento dal gioco per telefono, quello che chiede il catalogo nazionale completo, e la serie di cromatici legati al riempimento del catalogo di un titolo. Sono voci senza scadenza, perché le consegna il deposito e non la banca, e per questo non competono con il lavoro urgente; vanno però registrate perché sono collezionabili distinti nel senso adottato qui, e perché una di esse lega l'asse dei doni a quello delle sfide mentre le cromatiche lo legano al riempimento dei cataloghi.

Due altri assi hanno ricevuto in questa lettura la loro fonte. Gli esemplari sottolivellati, cioè ottenuti a un livello inferiore a quello in cui la specie evolve per natura, sono una proprietà che non si ricrea giocando in un altro titolo e che si perde con la via che la produceva; la guida dedicata \cite{sottolivellati} dichiara le dieci voci più difficili con la loro probabilità, e la loro utilità qui non sono le probabilità, che riguardano chi cattura e non chi compone, ma il fatto che dicono quali voci siano di fatto irraggiungibili per via di gioco entro la scadenza. Gli esemplari cromatici garantiti in gioco sono sette voci con il loro luogo, di cui tre nei titoli sotto scadenza; su questa seconda fonte va dichiarato un difetto nostro e non della fonte, perché l'estrattore che ne ha ridotto la pagina a testo ha perso i nomi delle specie conservando i luoghi, e il grezzo conservato accanto al derivato permette di recuperarli senza scaricare di nuovo.

\section{Gli scambi in gioco, contati sulla fonte di primo livello}

Le ultime due classi che il lavoro nominava senza contarle sono gli scambi in gioco e gli incontri che una condizione sblocca. Ciò che le unisce non è la rarità, che è tempo speso e non una proprietà del dato, ma il fatto che il dato porti un campo il cui valore, per quella specie, non si ottiene per altra via: in uno scambio è l'allenatore di provenienza con il suo identificativo e spesso un soprannome fissato dal gioco, in un incontro condizionato è la terna fra luogo, tipo di casella e specie.

Gli scambi sono stati contati sulle tabelle del verificatore, che è la fonte di primo livello, e non sulle due enumerazioni comunitarie che li contavano da fuori. Le voci di tabella sono duecentotrentotto e quelle distinte duecentotrentatre, su centocinquantadue specie. I due numeri differiscono perché le tabelle di una coppia di titoli si sovrappongono per costruzione, cioè esiste una tabella comune alla coppia e accanto una per versione, e la stessa voce compare in entrambe; la deduplicazione è dichiarata nel programma invece di essere lasciata implicita. Va notato per onestà che il totale della fonte comunitaria coincide con il nostro, e che la coincidenza non è ancora un accordo: le due enumerazioni non sono state confrontate riga per riga, e finché non lo saranno un numero uguale può nascere da insiemi diversi.

Il dato che decide la producibilità non è il totale ma la colonna dei valori di personalità scritti nella fonte. Dove la fonte lo scrive, cioè in diciannove voci di terza generazione, sedici di quarta, sette di quinta e quattro di ottava, l'esemplare è riproducibile byte per byte senza alcuna ricerca di semi, perché il valore che determina identità, sesso, abilità e cromaticità è un dato della tabella e non un'estrazione. Dove manca, perché il gioco lo genera al momento della consegna, la fedeltà torna indecidibile nello stesso senso già stabilito per i doni moderni. In prima e seconda generazione quel valore non esiste affatto, quindi la legittimità è banale e la fedeltà si esaurisce nei valori individuali, che la fonte scrive per tutte le voci di seconda generazione.

Due osservazioni valgono più del conto. La tabella dei doni di un'applicazione per console fissa sta nella fonte fra le tabelle di scambio e non fra i doni, quindi una classe a via chiusa trovata attraverso una fonte comunitaria risulta enumerabile dalla fonte di primo livello, con il valore di personalità di ciascuna voce. E quattro voci di terza generazione stanno nella tabella degli scambi di un gioco per console fissa, cioè richiedono un titolo che il progetto non possiede: la via esiste nel mondo e non per noi, che è una condizione diversa dalla via chiusa e va tenuta distinta.

\section{Gli incontri condizionati, e la correzione dell'esempio che li ha aperti}

La classe è nata da un esempio dell'utente, cioè un esemplare che si cattura su un'isola che compare a intermittenza e che conserverebbe quell'isola come luogo d'incontro. La verifica sulle tabelle corregge la premessa, e la correzione va scritta perché altrimenti sarebbe entrata nei documenti come un fatto.

L'isola non ha un luogo d'incontro proprio: nella fonte è l'area d'erba di una rotta, e quel medesimo luogo porta accanto le proprie aree d'acqua e di pesca con le loro specie ordinarie. Un esemplare catturato là registra dunque la rotta, non l'isola. Ciò che lo rende comunque irripetibile è un fatto diverso e più preciso, cioè che in quel luogo l'erba ospita quella specie sola e che quella specie non compare in nessun'altra area del titolo. La traccia sul dato è la terna, non il nome del luogo.

Da questa correzione discende un criterio meccanico, applicato a tutte le tabelle selvatiche dalla prima alla quinta generazione, cioè ventuno archivi. La prima prova è il tipo di casella, perché la fonte classifica ogni area e alcuni tipi sono essi stessi una condizione: sono divisi in due famiglie, la condizione di evento, come lo sciame annunciato, la grotta nascosta, la gara di cattura, l'albero del miele e le aree della zona safari, e la condizione di metodo, come le due mosse che aprono un incontro, che chiedono al giocatore un modo di cercare invece di un fatto avvenuto. Le specie che fra le sole vie selvatiche vengono da una condizione sono centotrentasette. La seconda prova è l'area monospecie esclusiva, ed è quella che coglie il caso dell'isola: quarantanove specie, e fra esse i casi che la comunità nomina da vent'anni. La terza è il luogo dedicato, cioè un luogo che ospita una sola specie contando tutte le sue aree: dodici specie, ed è il criterio più stretto perché il solo luogo scritto sul dato basta a stabilire la provenienza.

Due righe del quadro d'insieme saltano all'occhio e hanno una spiegazione che non è un difetto. Le ottanta specie delle riedizioni della seconda generazione vengono dalla zona safari, che in quei titoli è costruita dal giocatore disponendo blocchi e attendendo giorni, quindi ogni sua area porta un tipo proprio e ogni specie che vi si trovi soltanto risulta condizionata. Le quarantotto della seconda coppia di quinta generazione vengono dalle grotte nascoste, che sono il caso più puro della classe, perché la loro casella si ripopola secondo una condizione e alcune specie non esistono altrove nel gioco.

Il limite della misura va dichiarato perché ne cambia la lettura: le tabelle selvatiche non contengono incontri fissi, doni e scambi, quindi la frase \emph{unica via} vale sulle vie selvatiche e non su tutte. Il caso concreto è la specie dell'esempio, che nei tre titoli della sua generazione ha anche un uovo consegnato da un personaggio, cioè un incontro fisso e non un'area.

\section{Che cosa il criterio esclude, e la domanda che resta}

Applicato alla lettera alle due classi nuove, il criterio di produzione le esclude quasi tutte, e la ragione è la stessa per entrambe: le cartucce sono possedute, quindi la via esiste. Le eccezioni che il criterio riconosce subito sono tre e già enumerate, cioè gli scambi di un'applicazione dismessa, i quattro scambi dei giochi per console fissa e le voci delle riedizioni per console corrente, che scadenza non ne hanno affatto.

Resta però una tensione che il criterio non risolve. Alcune di queste vie esistono e sono praticamente impercorribili nei giorni che restano: l'isola compare in un giorno su decine di migliaia secondo il valore di personalità di un esemplare posseduto, le caselle in cui si trova una certa specie d'acqua sono sei su quasi cinquecento e si spostano, la zona safari delle riedizioni richiede giorni di attesa per ogni blocco disposto. La domanda se una via aperta e non percorribile entro la scadenza vada trattata come chiusa non è tecnica, ed è stata portata all'utente invece di essere decisa dentro un documento: la decisione che ne è seguita, cioè trattarla come chiusa e produrre tutto, è registrata fra le decisioni di questo lavoro e cambia il predicato del criterio da \emph{esiste una via} a \emph{esiste una via percorribile nel tempo residuo}.

\section{Il confronto con una enumerazione esterna, e la coda che ne è uscita}

L'ultimo strumento di questo capitolo è il confronto con il foglio comunitario del catalogo vivente, tenuto deliberatamente come controprova e non come autorità: le due enumerazioni non vengono fuse e nessuna delle due decide sull'altra, si mettono una accanto all'altra e le divergenze si classificano.

Il conto è di milletrecentottantanove voci per il foglio contro milletrecentosessantasette per noi, con i medesimi milleventicinque numeri di catalogo distinti e ottocentonovantadue specie su cui le due concordano esattamente. La differenza fra le due somme inganna se letta da sola, perché è la differenza fra due scarti che vanno in versi opposti e si compensano in parte: il foglio conta più di noi su centosette specie e noi più del foglio su ventisei. Il primo verso misura dove siamo ciechi, il secondo dove abbiamo rumore, e da esso sono venuti due difetti nostri, cioè le differenze di sesso che il campo della forma non separa e alcune posizioni di riempimento delle tabelle prese per forme.

Un secondo confronto della stessa specie è stato fatto sugli identificativi di allenatore documentati \cite{bulbapedia-id-notevoli}, che è una tabella di ottocentonovantasei righe utili con l'identificativo, il nome dell'allenatore e l'evento. Quattrocentosettantatre righe corrispondono alle nostre liste sulla coppia fra allenatore e identificativo; trecentouno corrispondono sul solo numero, e la corrispondenza è dichiarata debole e tenuta in una colonna a parte perché due eventi diversi possono condividere un numero; tre righe la fonte stessa le dichiara non appartenenti ad alcun esemplare da collezione; centodiciannove restano non coperte. Guardate una per una, quelle centodiciannove sono in grandissima parte distribuzioni coreane e giapponesi, cioè esattamente la classe che una decisione precedente aveva escluso per la barriera della lingua: la copertura del lavoro sulle distribuzioni documentate è dunque buona e il buco sta dove qualcuno ha deciso che stesse, il che è una proprietà molto diversa dall'avere un buco.

Da questi confronti nasce anche una coda di produzione di natura nuova, cioè le distribuzioni che non hanno mai lasciato una carta. Per ogni altra classe il generatore parte dal file che la distribuzione consegnava e ne compone i campi lasciati liberi; qui la carta non esiste e non è mai esistita, perché la consegna avveniva presso un apparecchio che scriveva direttamente nel salvataggio, oppure dentro il gioco, oppure attraverso il servizio in rete della quinta generazione, e il solo dato disponibile è quello che un archivio ha registrato mentre l'evento accadeva. Quella coda porta due casi in cui un campo non è un valore e che vanno decisi uno per uno: il campo che appartiene a chi riceve, che non si copia ma si scrive con l'allenatore del progetto, e il campo che dichiara un'alternativa invece di una scelta, tipicamente due abilità possibili o un livello che l'archivio non conosce.

\section{La terza enumerazione, che rompe la parità}

Due enumerazioni che divergono dicono che una delle due sbaglia e non quale, ed è la ragione per cui una terza vale più della somma delle prime due. La terza è quella di un tracciatore di catalogo vivente che pubblica i propri dati come JSON sotto licenza permissiva \cite{pokepc-dati}, il che la rende confrontabile riga per riga invece che a occhio: non è una lista di specie ma una disposizione in scatole, cioè l'elenco delle caselle che quello strumento ritiene si debbano riempire perché un deposito sia completo. È precisamente la domanda che la nostra lista di spunta dichiara indeterminata, cioè quali forme il deposito conti come casella separata, e su cui nessuna fonte di primo livello si pronuncia.

Il conto è di milletrecentottantasette voci contro le milletrecentosessantasette nostre e le milletrecentottantanove del foglio, con ottocentonovantatre specie su cui la terza enumerazione e la nostra concordano esattamente. Il fatto che conta non è lo scarto ma la sua forma: le due enumerazioni esterne, prodotte da autori diversi con metodi diversi, distano fra loro due voci su quasi millequattrocento, mentre la nostra dista da entrambe una ventina. Una convergenza di quel genere non dimostra che le due abbiano ragione, ma sposta l'onere della prova sulla terza, che siamo noi.

La classificazione delle divergenze è qui migliore che sul foglio, perché il dato porta per ogni voce i contrassegni che dicono che cosa essa sia. Delle trecentosessantadue voci che quella enumerazione conta oltre la specie base, centocinquantasette sono forme cosmetiche, centotré sono forme femminili e centotré sono forme che il campo della forma separa davvero. Le megaevoluzioni, le forme gigamax e quelle di sola battaglia non compaiono affatto, il che corrobora per via indipendente la scelta di tenerle fuori: uno strumento scritto da qualcun altro, per un altro scopo, esclude le stesse categorie che il nostro perimetro escludeva.

Una avvertenza di metodo va registrata perché ha rischiato di produrre un numero sbagliato con l'aria di essere esatto. Le disposizioni disponibili per il deposito sono sette, e la prima stesura del programma di confronto assumeva che differissero per il solo ordinamento; non è vero, perché due di esse danno una casella propria alla forma gigamax e due si dichiarano minime, cosicché i totali si raggruppano in millequattrocentoventicinque, milletrecentottantasette e milletrecentosettantatre. Il programma quindi non assume ma misura, dichiara quale disposizione usa e di quanto quella scelta muova il totale. Una cella, inoltre, non è sempre un identificativo: può essere un oggetto con un contrassegno, e identificarla con il solo identificativo faceva sparire trentotto caselle.

Due voci minori chiudono il quadro e valgono perché cadono su due questioni già decise. La prima è che quella enumerazione conta sessantatré configurazioni per la specie il cui dolcetto non è un campo della forma, numero che un foglio comunitario indipendente adotta a sua volta \cite{foglio-arca}, esattamente il numero che una decisione di questo lavoro aveva adottato contro il nove che la struttura del dato suggerisce. La seconda è che dichiara centotré specie con differenze di sesso contro le centodue che il progetto ha chiuso su due fonti, e la specie in più è una sola: la pagina enciclopedica non la nomina affatto, quindi il conto resta centodue e la voce in più è registrata come divergenza non corroborata invece che come correzione.

Due strumenti letti nello stesso giro chiudono il quadro sul perché la scelta del profilo non si possa rimandare all'infinito. Il primo genera la disposizione delle scatole a partire da più di quindici opzioni e tiene sopra di essa cinque contrassegni per esemplare, cioè l'allenatore proprio, la cattura nel titolo o nella regione di origine, la cromaticità, la sfera e la provenienza dall'applicazione per telefono \cite{monarium}. Il secondo è un foglio del 2022 che tiene insieme sei tracciatori distinti, cioè il catalogo vivente, gli esemplari di taglia massima, i marchi di origine, i marchi ordinari, i fiocchi e le sfere \cite{foglio-collezione-thundrosaur}. La coincidenza fra quei sei tracciatori e i sei assi che questo lavoro ha ricostruito per conto proprio, quattro anni dopo, è la conferma indipendente più forte che l'insieme degli assi sia quello e non un altro: due enumerazioni fatte senza conoscersi arrivano alla stessa partizione.
